\begin{helper}
Let \(H\) be a history cell, \(T\) a target event, \(I\) the fixed vertex set,
\(R\) the recorded-coordinate set, and \(U\) the terminal-coordinate set.
Assume there is a global red exceptional package: a number \(N_R^0\) with
\[
N_R^0\le 3\varepsilon_1 k^2
\]
such that for every branch \(\beta_0\in H\) there is a red exceptional set
\(E_R(\beta_0)\subseteq U\), with \(|E_R(\beta_0)|\le N_R^0\), which contains
every present red staged-open coordinate of any completion \(\omega\) that
has the same embedding as \(\beta_0\), agrees with \(\beta_0\) on recorded
coordinates, agrees with \(\beta_0\) on all blue terminal coordinates, is
regular, and has \(I\) independent in the final graph.

Assume also \(T\subseteq H\), every \(\omega\in T\) is regular, every
\(\omega\in T\) has \(I\) independent in its final graph, and the red
present staged-open support
\[
R(\omega)=
\{e\in \operatorname{StagedOpen}(\omega,\varepsilon,I):
  X_e(\omega)=1\text{ and }e\text{ is a red coordinate}\}
\]
has cardinality \(u_R\) for every \(\omega\in T\).

Then there are a natural number \(N_R\), an abstract finite label type
\(\mathcal J\), and a branch-indexed decoding function
\[
\operatorname{redSupport}:H\times
\mathcal J\to \{\hbox{finite coordinate sets}\}
\]
such that
\[
N_R\le 3\varepsilon_1 k^2,\qquad
|\mathcal J|\le \binom{N_R}{u_R}.
\]
For every target sample \(\omega\in T\) and every compatible branch
\(\beta_0\in H\), there is a label \(J\in\mathcal J\) for which
\[
\operatorname{redSupport}(\beta_0,J)=R(\omega).
\]
For this assigned label, \(\operatorname{redSupport}(\beta_0,J)\) has
cardinality \(u_R\), is contained in \(U\), and consists only of red
coordinates.
\end{helper}
\begin{proof}
Choose \(N_R=N_R^0\).  For each branch \(\beta_0\in H\), the exceptional
package gives a set \(E_R(\beta_0)\subseteq U\) of size at most \(N_R\).
Pad this set by dummy slots until it has \(N_R\) slots.  The dummy slots are
bookkeeping positions, not terminal coordinates.  Let the abstract label type
\(\mathcal J\) be the family of \(u_R\)-subsets of the \(N_R\) padded slots.
Thus
\[
|\mathcal J|\le\binom{N_R}{u_R}.
\]

For a branch \(\beta_0\), define \(\operatorname{redSupport}(\beta_0,J)\)
by taking from the package \(E_R(\beta_0)\) the actual red terminal
coordinates occupying the slots selected by \(J\); dummy slots decode to no
terminal coordinate.  This definition uses the branch \(\beta_0\) itself,
which is the datum that determines the exceptional package, not merely the
branch's blue terminal trace.
Thus the decoder is genuinely branch-indexed.  If two branches have the same
blue terminal values but different embedding or recorded data, the global
red exceptional witness may supply different exceptional sets for them; the
same abstract label \(J\) is decoded separately through the appropriate
set \(E_R(\beta_0)\).  This is exactly the input needed later by the
branch-label construction: after a representative branch is chosen, \(J\) is
interpreted through that representative's own exceptional package.

Now fix a target sample \(\omega\in T\) and a branch \(\beta_0\in H\) which
is compatible with \(\omega\) in the sense required in the conclusion: the
two samples have the same embedding, agree on the recorded coordinates, and
agree on all blue terminal coordinates.  By the target hypotheses,
\(\omega\) is regular and has
\(I\) independent in the final graph.  By compatibility, it has the same
embedding as the branch, agrees on the recorded coordinates, and agrees on
all blue terminal coordinates.  The exceptional package for \(\beta_0\)
therefore contains every element of the red present staged-open support
\(R(\omega)\) in \(E_R(\beta_0)\).  Since \(|R(\omega)|=u_R\), choose the
label \(J\) whose padded slots are exactly the slots occupied by
\(R(\omega)\) inside the padded list for this same branch \(\beta_0\).  The
decoding for this \(J\) is precisely \(R(\omega)\), because each real slot
selected by \(J\) is an element of \(R(\omega)\) and every element of
\(R(\omega)\) has been selected; no dummy slot is needed for this assigned
label.

The final three properties are only asserted for this assigned label, exactly
as in the Lean statement.  They follow from the fixed-count hypothesis
\(|R(\omega)|=u_R\), from \(E_R(\beta_0)\subseteq U\), and from the definition
of \(R(\omega)\) as the red-coordinate part of the present staged-open
support.
\end{proof}
